\subsubsection{機能指向・構造化設計}

\term{機能指向設計}{Function-Oriented Design} または\term{構造化設計}{Structured Design} は、ソフトウェアの主要機能を見つけ、それを上位から下位へ分解する古典的な方法である。

この方法では、データフロー図や処理記述を使い、入力がどの処理を通って出力へ変換されるかを考える。業務処理の流れが明確で、機能分解しやすいシステムで理解しやすい。一方、状態やデータ構造、イベント、オブジェクト間の相互作用が複雑な場合は、別の方法と組み合わせる必要がある。
